\begin{theorem}\label{night:P3K.34}
\StatementRecord{night:P3K.34}{false as stated; correction proved}{S6\_06\_BLIND: GPT-6 Astra (OpenAI, ChatGPT Pro mode, effort Pro), recorded re-derivation in the checkpoint dated 18 November 2026}{\cref{night:P3F.26}, \cref{night:P3F.29}, \cref{night:P3F.32}, \cref{night:P3K.65}, \cref{night:P3K.68}, \cref{night:P3K.76}, \cref{night:P3K.77}}
\par\noindent\textbf{Audit correction, 17 November 2026.} The universal OS49 fan formula below is false. For \(k=1\), the original negative-diagonal fan has \(S=1\) in its primitive \(H/K\) marking, so that formula gives \(0\); P3K.65 proves the actual positive \(r\)-loop class \(-3\). Its shorter negative-fan class is \(-k\), while the positive-fan class is \(0\). The old statement below is retained as the refuted record. The proof following it is the corrected local page and framing calculation, with physical pair \((2m,8m)\), and does not prove the old signs or the old all-fan formula. The original argument is preserved in \path{records/attempts/post19_recovery/oc_checkpoint19/notes/part3_os49_full_framing.tex}. The current vertical and all-fan formulas are P3K.76 and P3K.77.
\medskip


Retain the actual component-supported OS49 families of P3F.31--32,
with \(\lambda=r^{3k}\), \(k>0\), and cap vector
\(n(1,-1,0,0)^t+d\delta\), where \(n=\pm1\) and \(d\in\mathbb Z\).
For one positive physical \(r\)-loop, in the initial positive-diagonal
Cox product the individual integers and relative coordinates are
\[
 k_{\rm Jin}=-6k,\qquad \ell_{\rm Jin}=-24k,\qquad
 (m_{\rm rel},n_{\rm rel},s)=(6k,0,0).
 \tag{OC.1}
\]
The pair is the same for both norm signs. The line used has its actual
\(I_2\) multidegree \((-1,+1)\) and relative square \(-2k\).
For every actual allowed index-one fan \(T(v,w)\), write the physical
delta cocharacter as \(\delta=a v+b w\), and set \(S=a^2-ab+b^2\).
Its oriented return class is
\[
 [\Sigma_{\Phi,T}]=\frac{3k}{4}(S-1)(S-3)\pmod{28}.
 \tag{OC.2}
\]
At fixed \(k\), its exact range is
\(k\{0,4,6,8,14,18,20,22\}\pmod{28}\).
These conclusions retain the component line, both signed norm maps,
the simultaneous closed Cox product, and the actual common-exterior
frame comparison. They use the stated physical parameter loop.

\end{theorem}
\begin{proof}
\begin{enumerate}[label=\arabic*.]
\item \par

\medskip\noindent\textbf{1. The actual one-cap specialization and both solid-torus maps}

\par\nobreak

\item Retain the actual source, section and component line
\[
y^2=x^3-3t^3(t-1)x+2t^4(t-1)(t-\tfrac12),\quad
G=(t^2,t^2),\quad\mathcal M=\mathcal O(G-O),\quad Q=3G.
\]
The finite \(IV^*\) cap at \(t=0\) has order \(m=3k\),
full vector \(\nu\ell_0+d\delta\), \(\nu=\pm1\), and the literal
Cartier/norm maps of F26/F29/F32. The selected \(I_2\) cusp is at
\(t=\infty\); the original type-II fibre at \(t=1\) is retained.
The actual scalar family is \(\lambda=r^m\). Its endpoint-to-initial
ordinary closing is \(t^\nu\).

\item Here is the complete specialization morphism from RP.1-\/-7. Set
\[
h=1,\quad l=0,\quad n_A=\nu,\quad n_B=0,\quad
 d_A=d,\quad d_B=0,\quad M=m,\quad\epsilon=\nu.
\]
At the second site the root cover is the identity, its root deck and
norm action are the identity, and its ordinary-to-cap logarithmic
cochain is zero. Thus its quotient is exactly the original type-II
source/Cartier chart. In the coordinate \(q=t-1\), the source section
is \(((1+q)^2,(1+q)^2)\). At \(q=0\) it is the smooth point
\((1,1)\) of \(y^2=x^3\), distinct from \(O\); its coordinate
\(\tau=-x/y\) is \(-1\). Therefore the \(l=0\) terminal Cartier
ratio is \((\tau+1)/\tau\), the ratio for the original horizontal
divisor \(G-O\), with no exceptional graph blowup and no extra
vertical divisor. This is the actual source and line arrow at this
site. At the first site every root, Cartier and finite-character map
is unchanged. All the following maps are therefore maps of the
one-cap family itself.

\item Let \(y_B\) denote the coefficient of the actual contracting period
\(B\). On the first root \(q_A=z^m\), write
\(a=\arg z/(2\pi)\), \(\alpha_A=ma\). The plus-log arrow is
\[
y_{\rm cap}=y_B+\nu a,\qquad
(a,y_{\rm cap})\longmapsto(a+1/m,y_{\rm cap}+\nu/m).
\]
Its invariant cap circle is \(u_B=my_B+\nu\alpha_A\). The filled
ordinary meridian is \(m\alpha_A-\nu\gamma_B\). Define the exact
outward coordinates \(\gamma_{\rm out}=-y_B\),
\(u_{\rm out}=-u_B\). Then
\(u_{\rm out}=m\gamma_{\rm out}-\nu\alpha_A\), and the root cut
changes \(u_{\rm out}\) by \(-\nu\). The positive weighted Hopf
angle is consequently \(\beta=-\nu y_B\).

\item Put \(T=t/(t-1)\). On the ordinary annulus the phases are
\[
\phi_1=\nu(m\gamma_{\rm out}-\nu\arg T/(2\pi)),\qquad
\phi_2=\nu\gamma_{\rm out},\qquad
m\phi_2-\phi_1=\arg T/(2\pi).
\]
Choose the unique positive radii with
\(r_1^2+r_2^2=1\) and \(r_2^m/r_1=|T|\); then
\((z_1,z_2)=(r_1e^{2\pi i\phi_1},r_2e^{2\pi i\phi_2})\)
has \(T=z_2^m/z_1\). The extending maps at the two solid tori are
\[
\begin{aligned}
T^{-1}=z:\quad &(z,u_\infty)\longmapsto
(zR^m e^{2\pi i\nu m u_\infty},R e^{2\pi i\nu u_\infty}),
&&R^2+|z|^2R^{2m}=1,\\
T=z^m:\quad &(z,u_0)\longmapsto
(R e^{2\pi i\nu u_0},zR^{1/m}e^{2\pi i\nu u_0/m}),
&&R^2+|z|^2R^{2/m}=1.
\end{aligned}
\]
The root cut changes their phases by integral amounts, using
\(\nu^2=1\). Each inverse first recovers the positive radius from
the nonzero coordinate, then an argument lift and the remaining
root coordinate. Different lifts differ precisely by the actual
root cut. The maps have positive radial derivative and agree with
the ordinary overlap after the specified analytic coordinate units
are retained. On the actual source/line sheet a stabilizer satisfies
\(\nu a+mb=0\), hence \((a,b)=(mj,-\nu j)\); its action is
\(A^{mj}B^{-\nu j}=1\) by the literal norm. Thus the sheet-zero
map introduces no residual quotient.

\item For a regular \(\tau\notin\{0,1,\infty\}\), the page is
\[
\eta^m-\tau\xi=1,\qquad T=\tau+\xi^{-1}.
\]
The positive weighted radial map is
\((\xi,\eta)\mapsto(\rho^m\xi,\rho\eta)\), with
\(\rho^{2m}|\xi|^2+\rho^2|\eta|^2=1\). This is the actual complex
page of the polynomial \(z_2^m-\tau z_1\). Its degree over the
source base is \(m\); it has one point of ramification index \(m\)
over the finite root and the binding, and \(m\) unramified points
over the retained type-II fibre. Its compact base has genus zero.

\par

\medskip\noindent\textbf{2. The plus-origin square and the mixed boundary map}

\par\nobreak

\item Keep F32\textquotesingle s actual periods and prescribed unit logarithms:
\[
u=(2s,s)^t+U(q),\quad
\gamma_B=(-3s,-s+mv)^t+C(q),\quad\mathcal D=3U+2C,
\]
\[
H=-u-\gamma_B,\qquad
K=3u+2\gamma_B=\mathcal D+(s+2mv)\delta,
\qquad s=\frac{\log q}{2\pi i},\quad v=\frac{\log r}{2\pi i}.
\tag{OC.3}
\]
The valuation matrix has columns \((2,1),(-3,-1)\) and determinant
one. The two actual products have PLUS origins
\[
\xi_{\rm ord}=\zeta+o,\quad\xi_{\rm cusp}=\zeta+\sigma,\quad
 o=-\frac{\nu s\mathcal D}{2m}-\frac{\nu s^2}{4m}\delta-\nu sv\delta,
\quad\sigma=\nu v\mathcal D+\nu mv(v-1)\delta.
\tag{OC.4}
\]
Thus cusp-to-ordinary is \(o-\sigma\). The ordinary closing logarithm
is \(c=-\nu s\delta\), whereas
\(\sigma(v+1)-\sigma(v)=c+\nu K\). The bounded FORWARD endpoint
is therefore \(-\nu K\). Its full angular path is
\[
(Z,V)\longmapsto
\left(U_K^{-\nu\theta}Z,
(W_K\lambda_0^2)^{-\nu\theta}
 e^{-2\pi i\nu m\theta(\theta-1)}V\right).
\tag{OC.5}
\]
The source-root \(s\mapsto s+2m\) difference retains its integral
boundary gauge. All powers use the actual logarithms of the bounded
units. The quadratic scalar factor is the global physical-circle
map on the source, line and finite cap, with its endpoint retained.
For full multiplier \((A,D)\), its local Cox lift is
\[
(z_0,z_4,z_1,z_2)\longmapsto
((AD)^{-1}z_0,z_4,Az_1,Dz_2).
\tag{OC.6}
\]
It preserves \(q=z_0z_4z_1z_2\) and the actual period and cap arrows.
Here it is used for the positive fixed-star product and its carried
normal references; no simultaneous marked product for all curve
neighbourhoods is asserted.

\item Let \((x,y_B)\) be point angles of \((u,\gamma_B)\) and
\((h_0,\kappa)\) those of \((H,K)\). Their literal maps are
\[
\binom{h_0}{\kappa}=\begin{pmatrix}2&-3\\1&-1\end{pmatrix}
\binom{x}{y_B},\qquad
\binom{x}{y_B}=\begin{pmatrix}-1&3\\-1&2\end{pmatrix}
\binom{h_0}{\kappa}.
\tag{OC.7}
\]
The incoming compression graph \(\kappa=g(h_0,w)\) moves by
\(-\nu\theta\), by OC.5. Hence
\[
\theta=-\nu(\kappa-g),\qquad\beta=-\nu y_B,
\tag{OC.8}
\]
\[
\binom{h_0}{\theta}=\begin{pmatrix}2&3\nu\\-\nu&-1\end{pmatrix}
\binom{x}{\beta}+\binom0{\nu g},\qquad\det=1.
\tag{OC.9}
\]
Both signs have the same positive degree. The source cycle \(w\) is
unchanged. The Tate lengths of \(u,H\) have positive leading terms
\(2\operatorname{Im}s,\operatorname{Im}s\), so their complex source
orientations agree on the chosen collar. Consequently
\[
\partial V_m\longrightarrow E_H\times S^1_\theta
\quad\text{has degree }+1.
\tag{OC.10}
\]
Orient the compression by \(\partial W=F-E_H\). Its product trace
has boundary \(F\times S^1_\theta\). The following actual Hopf
coordinate fixes the normal orientation; its positive meridian
line is \(U=\mathcal M^*\). The physical line is carried by K68\textquotesingle s
explicit dual arrow at the end of the computation.

\par

\medskip\noindent\textbf{3. The full Cox normal frame and its signed reference}

\par\nobreak

\item The positive fan has charts
\[
(A,B,Z)=(q/V,V/Z,Z),\qquad
(A',B',Z')=(ZA,ZB,Z^{-1}),\qquad q=ZAB,
\tag{OC.11}
\]
with physical weights \((-1,+1,0)\). Put
\(y=AB\), \(b=(|B|^2-|A|^2)/2\). The positive \(K\) height
\(\chi_K\) decreases with \(\log|V|\), so
\(\partial_b\chi_K<0\). Retain this reflection at both references.
The global quaternion is
\[
q_H=(1+jZ)(A+j\overline B)
       =(Z^{-1}+j)(A'+j\overline{B'}).
\tag{OC.12}
\]
Physical \(\delta\) acts by \(q_H\mapsto q_He^{-i\phi}\).
At \(Z=0\),
\[
q_Hi\overline{q_H}=-2bi+2\operatorname{Im}(y)j
                               -2\operatorname{Re}(y)k.
\tag{OC.13}
\]
Its determinant relative to \((\operatorname{Re}y,\operatorname{Im}y,b)\)
is \(-8\); the circle-first quotient determinant at \((A,B)=(1,0)\)
is \(-1\). Together with OC.10 these give the positive normal
Hopf meridian. The initial comparison in
\((\operatorname{Re}y,\operatorname{Im}y,-b)\) is positive triangular.
Its shear and positive diagonal contract. For the signed coordinate
\(\beta\), apply
\[
T_{-\nu}=\operatorname{diag}(1,1,-\nu)
\tag{OC.14}
\]
to BOTH references. The residual comparison is the same positive
triangular matrix; conjugation by \((-\nu)T_{-\nu}\in SO(3)\)
adds no free framing class.

\item For a bounded complex multiplier \((Z,V)\mapsto(aZ,b_1V)\), the
actual chart map is
\((Z,A,B)\mapsto(aZ,b_1^{-1}A,(b_1/a)B)\). The quaternion formula
uses its UNITARY phases. With \(p=b_1^{-1}\), \(q_1=b_1/a\), choose
\(L^2=pq_1\), \(R^2=p/q_1\), \(LR=p\) on the actual interval;
then \(q_H\mapsto Lq_HR\). Carry the positive moduli by
\[
a_t=(a/|a|)|a|^t,\quad b_t=(b_1/|b_1|)|b_1|^t,\qquad
(Z,A,B)\mapsto(a_tZ,b_t^{-1}A,(b_t/a_t)B),\quad0\le t\le1.
\tag{OC.15}
\]
This preserves \(q\), both chart transitions and both normal
references. At \(t=0\) the quaternion identity is literal. Higher
\(q\)-unit terms have zero normal linear part because \(q=ZAB\)
is quadratic in the normal directions. The closed rotation is
constant along \(F\), extends over its \(D^3\), and retains its
actual Spin endpoint.

\par

\medskip\noindent\textbf{4. The binding phase and the complete normal clutch}

\par\nobreak

\item The actual regular binding coordinate is
\(q_b=(T-\tau)(T/\tau)^{-1}\), with derivative one at \(T=\tau\).
The polynomial phase gives \(\alpha_b=\phi-m\beta\). By OC.9,
\[
\beta=2\theta+\nu(h_0-2g),\qquad
\alpha_b=\phi-2m\theta-\nu m(h_0-2g).
\tag{OC.16}
\]
The source circle term is retained. In the quotient of the collar
normal tangent \((\cos\alpha,\sin\alpha,\rho,0)\), \(\rho>0\),
the actual map sends
\((v_q,v_r,v_\eta)\) to
\((v_q-v_r(\cos\alpha,\sin\alpha)/\rho,v_\eta)\).
The two Cartesian parameter vectors give
\[
G(\alpha,\theta)=\begin{pmatrix}
-\sin\alpha&\cos\theta\cos\alpha&\sin\theta\cos\alpha\\
\cos\alpha&\cos\theta\sin\alpha&\sin\theta\sin\alpha\\
0&\sin\theta&-\cos\theta
\end{pmatrix}\in SO(3).
\tag{OC.17}
\]
A compression frame transported by the actual \(K\) flow differs
from the Cox frame by \(R(p)\), independent of \(\theta\). Its
\(S^2\) boundary map extends over \(D^3\). The compression is a
solid torus minus a ball, and its oriented normal three-bundle has
zero \(w_2\) in the actual stably parallelizable orbit \(S^5\).
The complete quotient matrix obeys the literal factorization
\[
G(\alpha,\theta)=R_z(\alpha)G(0,\theta),
\tag{OC.18}
\]
so the full clutch is
\[
[R(p)^{-1}R_z(-\nu m(h_0-2g))]
[R_z(\phi-2m\theta)G(0,\theta)].
\tag{OC.19}
\]
The page is the fixed polynomial level \(\eta^m-\tau\xi=1\), so its phase \(\phi\) is fixed throughout this clutch. The first factor is pulled from the actual source two-torus; the
second is pulled from the parameter circle. Expansion of the
Maurer-\/-Cartan cubic for their product gives the sum of factor
cubics plus an exact form. Both factor cubics vanish by dimension.
With the carried collar connection this proves
\[
\langle p_1(\overline\nu),[\overline X]\rangle=0.
\tag{OC.20}
\]
All one-dimensional normal data and both signed references remain
in this calculation.

\par

\medskip\noindent\textbf{5. The actual component line, cap markings and physical pair}

\par\nobreak

\item The sheet-zero morphism above carries the exact root line to
\(N=\mathcal O(G_{\rm page}-O)\). At the unmodified type-II points
\(G\) and \(O\) stay distinct, so their intersection contribution
is zero. The source page has \(m\) unchanged \(I_2\) fibres,
\(m\) unchanged type-II fibres and one good root. Thus
\[
\chi(\mathcal O_{S_{\rm page}})=m/3=k,\quad
\sigma(S_{\rm page})=-8k,\quad N^2=-2k,
\quad p_1(TS_{\rm page})=-24k.
\tag{OC.21}
\]
The relative period torsor is trivialized with its exact binding
collar: a lifted type-II meridian has elliptic monodromy with
\(B_{\rm ell}-I=\left(\begin{smallmatrix}0&1\\-1&-1\end{smallmatrix}\right)\),
which is unimodular. Hence the integral monodromy coinvariants
vanish. The component terms remain finite skyscraper sheaves.
The rigidified fibrewise Picard/Poincare Hom map carries the line
under the resulting translations, retaining every component degree
and finite character.

\item The literal line subfan has rays \((j,\lceil j/2\rceil,1)\) and
component degrees \((-1,+1)\). Strict contraction gives \(m\)
ordered, separated norm shells; the positive-iterate height defects
prevent reentry after the first zero. The exact radial comparison
\(v\mapsto c(q)^{F(s)-s}v\), \(F(s+1)=F(s)+1\), is periodic,
commutes with the contracting lift and physical circle, and fixes
the source. It carries the actual line into the compact page.

\item The graph-normal pushoff uses these same annular/binding charts.
Outside the regular bands its radial homotopy misses the compression.
Inside each band the exact equation is
\(\rho=-\eta_0\cot\alpha\), \(\sin\alpha>0\), \(\cos\alpha\le0\).
The slightly enlarged boundary arcs have disjoint gaps, and the
inner semicircle includes into its half-disc by a homotopy
equivalence. After the source-circle product, each relative
\(H^2\) is zero, so restriction is injective. The normal meridian
line is therefore the carried page line throughout the comparison.
The fixed \(d\) character is retained on every sheet; its full
product is one. Its compact extension changes only an integral
fibre/base class with zero square and zero pairing with \(N\).

\item The incoming basis is \((u,\gamma_B)=(2H+K,-3H-K)\), determinant
one. The compression cap \((D_w^2\times S_H^1)\times S_K^1\)
has \(H^2\) injecting into the boundary and a flat actual line.
The regular cap \(T^2_{u,w}\times D^2_{\gamma_B}\) also has zero
line square and signature. They retain their different primitive
killed cycles. Their actual line restrictions and Novikov
additivity give \(b_U^2=4N^2=-8k\) and
\(p_1(T\overline X)=-24k\).

\item With the positive Hopf meridian \(U=\mathcal M^*\), the
sphere-bundle pushforward gives
\[
(k_U,\ell_U)=(-6k,-24k).
\]
K68\textquotesingle s actual dual/spinor arrow
\(Q(a,b)=(\overline b,-\overline a)\), together with the inverse
closing in OC.5, gives
\(\lambda_{\mathcal M}(F)=-\Sigma(F^{-1})=\Sigma(F)\).
It reverses both characteristic integrals, so the physical pair is
\[
(k_{\mathcal M},\ell_{\mathcal M})=(6k,24k)=(2m,8m),\qquad
A_{\mathcal M}=\frac{4k_{\mathcal M}-\ell_{\mathcal M}}{24}=0.
\tag{OC.22}
\]
Both cap and origin arrow reversals were used. The one-cap proof
therefore retains the physical positive-fan reference used in
K76/K77. The actual negative-fan/shorter comparison is K65; the
all-fan computation is K77. No old OC all-fan polynomial or
simultaneous marked-plumbing premise is reinstated.

\end{enumerate}
\end{proof}